\documentclass[12pt]{article}
% Shared LaTeX preamble for portfolio papers
\usepackage[margin=1in]{geometry}
\usepackage[T1]{fontenc}
\usepackage[utf8]{inputenc}
\usepackage{newtxtext}
\usepackage{microtype}
\usepackage{setspace}
\usepackage{xurl}
\usepackage{hyperref}
\usepackage{fancyhdr}
\usepackage{titlesec}
\usepackage{enumitem}
\usepackage{longtable}
\usepackage{array}

\hypersetup{
    hidelinks,
    colorlinks=false,
    pdftitle={ED400 Comprehensive Teaching Portfolio},
    pdfauthor={Piter Garc\'{i}a}
}

\doublespacing
\setlength{\parindent}{0.5in}
\setlength{\parskip}{0pt}
\setlength{\emergencystretch}{3em}

\pagestyle{fancy}
\fancyhf{}
\fancyhead[R]{\thepage}
\renewcommand{\headrulewidth}{0pt}
\setlength{\headheight}{15pt}

\newcommand{\PortfolioAuthor}{Piter Garc\'{i}a}
\newcommand{\PortfolioAffiliation}{University of Rochester, Warner School of Education}
\newcommand{\PortfolioProgram}{Initial Professional Teaching Certification --- K--12 Computer Science}
\newcommand{\PortfolioDate}{May 2026}

\newcommand{\apaheading}[1]{\begin{center}\textbf{#1}\end{center}}
\newcommand{\reference}[1]{\noindent\hangindent=0.5in\hangafter=1 #1\par}

\title{ED400 Comprehensive Teaching Portfolio}
\author{\PortfolioAuthor}
\date{\PortfolioDate}
\begin{document}
\maketitle

\section*{Portfolio Overview}
This single-file document serves as the umbrella entry point for the comprehensive teaching portfolio. The full portfolio is organized as eleven short papers: one personal statement and ten proficiency papers aligned to the program standards. For clean repository use, the substantive papers are stored individually in the \texttt{personal-statement/}, \texttt{proficiencies/}, and \texttt{tex/} directories.

\section*{Included Papers}
\begin{enumerate}[leftmargin=1.5em]
    \item Personal Statement
    \item Proficiency 1: Learner Development
    \item Proficiency 2: Learner Diversity and Equity
    \item Proficiency 3: Learning Environments
    \item Proficiency 4: Content Knowledge and Application
    \item Proficiency 5: Assessment of Learning
    \item Proficiency 6: Planning for Instruction
    \item Proficiency 7: Pedagogy
    \item Proficiency 8: Communication and Collaboration
    \item Proficiency 9: Reflection and Professional Growth
    \item Proficiency 10: Advocacy and Ethical Leadership
\end{enumerate}

\section*{Repository Guide}
\begin{itemize}[leftmargin=1.5em]
    \item \texttt{README.md} --- portfolio overview and quick-start guide
    \item \texttt{NAVIGATION.md} --- proficiency-to-artifact cross-reference map
    \item \texttt{artifacts/ARTIFACTS-REGISTRY.md} --- master artifact index with clickable links
    \item \texttt{personal-statement/} --- Markdown version of personal statement
    \item \texttt{proficiencies/} --- Markdown versions of P1--P10
    \item \texttt{tex/} --- LaTeX versions of portfolio materials
\end{itemize}

\section*{Compilation Note}
At this stage, this repository contains a strong organizational scaffold, a completed personal statement, artifact documentation, and concise draft proficiency papers in Markdown. The next step is to expand each proficiency paper to full rubric length (typically 8--15 pages with embedded artifact discussion) and mirror those final versions in \LaTeX{}.

\end{document}
